
\textbf{II. PHẦN TỰ LUẬN}
\begin{ex}
	Rút gọn: $B = \left( \frac{\sqrt{x}}{x-2\sqrt{x}} + \frac{x+2\sqrt{x}}{x-4} \right) : \left( 1 + \frac{1}{\sqrt{x}} \right)$; $x>0, x \neq 4$.
	\loigiai{
		Điều kiện: $x>0, x \neq 4$.
		Ta có:
		\begin{align*} B &= \left( \frac{\sqrt{x}}{\sqrt{x}(\sqrt{x}-2)} + \frac{\sqrt{x}(\sqrt{x}+2)}{(\sqrt{x}-2)(\sqrt{x}+2)} \right) : \left( \frac{\sqrt{x}+1}{\sqrt{x}} \right) \\ &= \left( \frac{1}{\sqrt{x}-2} + \frac{\sqrt{x}}{\sqrt{x}-2} \right) \cdot \frac{\sqrt{x}}{\sqrt{x}+1} \\ &= \left( \frac{1+\sqrt{x}}{\sqrt{x}-2} \right) \cdot \frac{\sqrt{x}}{\sqrt{x}+1} \end{align*}
		Vì $x>0$, nên $\sqrt{x}+1 > 0$.
		\begin{align*} B &= \frac{\cancel{(\sqrt{x}+1)}}{\sqrt{x}-2} \cdot \frac{\sqrt{x}}{\cancel{(\sqrt{x}+1)}} \\ &= \frac{\sqrt{x}}{\sqrt{x}-2} \end{align*}
		Vậy $B = \frac{\sqrt{x}}{\sqrt{x}-2}$.
	}
\end{ex}

\begin{ex}
	Cho phương trình: $x^2-2x-1=0$ không giải phương trình, hãy tính giá trị biểu thức sau: $T=(x_1+1)^2+(x_2+1)^2$.
	\loigiai{
		Phương trình $x^2-2x-1=0$ có $a=1, b=-2, c=-1$.
		$\Delta' = b'^2 - ac = (-1)^2 - 1(-1) = 1+1=2 > 0$.
		Phương trình có hai nghiệm phân biệt $x_1, x_2$.
		Theo hệ thức Vi-ét, ta có:
		$$ \begin{cases} x_1 + x_2 = -\frac{b}{a} = -\frac{-2}{1} = 2 \\ x_1 x_2 = \frac{c}{a} = \frac{-1}{1} = -1 \end{cases} $$
		Ta biến đổi biểu thức T:
		\begin{align*} T &= (x_1+1)^2 + (x_2+1)^2 \\ &= (x_1^2 + 2x_1 + 1) + (x_2^2 + 2x_2 + 1) \\ &= (x_1^2 + x_2^2) + 2(x_1+x_2) + 2 \\ &= ((x_1+x_2)^2 - 2x_1x_2) + 2(x_1+x_2) + 2 \end{align*}
		Thay giá trị từ hệ thức Vi-ét vào biểu thức T:
		\begin{align*} T &= ((2)^2 - 2(-1)) + 2(2) + 2 \\ &= (4 + 2) + 4 + 2 \\ &= 6 + 4 + 2 \\ &= 12 \end{align*}
		Vậy $T=12$.
	}
\end{ex}

\begin{ex}
	Giải hệ phương trình:
	$$ \begin{cases} 2x+3y=7 \\ 3x-y=5 \end{cases} $$
	\loigiai{
		Ta có hệ phương trình:
		$$ \begin{cases} 2x+3y=7 \quad &(1) \\ 3x-y=5 \quad &(2) \end{cases} $$
		Từ phương trình (2), suy ra $y = 3x-5$.
		Thế $y = 3x-5$ vào phương trình (1):
		\begin{align*}
			2x + 3(3x-5) &= 7 \\
			2x + 9x - 15 &= 7 \\
			11x &= 7 + 15 \\
			11x &= 22 \\
			x &= \frac{22}{11} \\
			x &= 2.
		\end{align*}
		Thay $x=2$ vào $y = 3x-5$:
		\begin{align*}
			y &= 3(2) - 5 \\
			&= 6 - 5 \\
			&= 1.
		\end{align*}
		Vậy nghiệm của hệ phương trình là $(x; y) = (2; 1)$.
	}
\end{ex}

\begin{ex}
	Trong một hộp thẻ được đánh số từ 1 đến 6. Lấy ngẫu nhiên 2 thẻ. Tính xác suất để các thẻ lấy ra có tổng các số trên thẻ bé hơn 5.
	\loigiai{
		Các thẻ được đánh số: $\{1; 2; 3; 4; 5; 6\}$.
		Lấy ngẫu nhiên đồng thời 2 thẻ. Các cặp thẻ có thể lấy (không gian mẫu $\Omega$):
		$\Omega = \{ (1;2); (1;3); (1;4); (1;5); (1;6); (2;3); (2;4); (2;5); (2;6); (3;4); (3;5); (3;6); (4;5); (4;6); (5;6) \}$.
		Số phần tử của không gian mẫu là $n(\Omega) = 15$.
		
		Gọi A là biến cố: “Tổng các số ghi trên 2 thẻ lấy ra bé hơn 5”.
		Các cặp thẻ có tổng các số bé hơn 5:
		\begin{itemize}
			\item Thẻ thứ nhất là 1:
			\begin{itemize}
				\item (1;2) (Tổng = 3 < 5)
				\item (1;3) (Tổng = 4 < 5)
			\end{itemize}
			\item Thẻ thứ nhất là 2 (và thẻ thứ hai lớn hơn 2 để không trùng lặp):
			% Không có cặp nào thỏa mãn vì (2;3) có tổng là 5.
		\end{itemize}
		Vậy, các kết quả thuận lợi cho biến cố A là:
		$A = \{ (1;2); (1;3) \}$.
		Số phần tử của biến cố A là $n(A) = 2$.
		
		Xác suất của biến cố A là:
		$P(A) = \frac{n(A)}{n(\Omega)} = \frac{2}{15}$.
	}
\end{ex}
\begin{ex}
	Đầu năm học, Hội khuyến học của một tỉnh tặng cho trường A tổng số 245 quyển sách gồm sách Toán và sách Ngữ văn. Nhà trường đã dùng $\frac{1}{2}$ số sách Toán và $\frac{2}{3}$ số sách Ngữ văn đó để phát cho các bạn học sinh có hoàn cảnh khó khăn. Biết rằng mỗi bạn nhận được một quyển sách Toán và một quyển sách Ngữ văn. Hỏi Hội khuyến học tỉnh đã tặng cho trường A mỗi loại sách bao nhiêu quyển?
	\loigiai{
		Gọi số sách Toán mà Hội khuyến học tỉnh đã tặng cho trường A là $x$ (quyển).
		Gọi số sách Ngữ văn mà Hội khuyến học tỉnh đã tặng cho trường A là $y$ (quyển).
		Điều kiện: $x, y \in \mathbb{N}^*$.
		
		Tổng số sách hai loại là 245 quyển, nên ta có phương trình:
		$x + y = 245 \quad (1)$
		
		Số sách Toán đã dùng để phát là: $\frac{1}{2}x$ (quyển).
		Số sách Ngữ văn đã dùng để phát là: $\frac{2}{3}y$ (quyển).
		
		Vì mỗi bạn nhận được một quyển sách Toán và một quyển sách Ngữ văn, nên số sách Toán đã phát bằng số sách Ngữ văn đã phát:
		\begin{align*} \frac{1}{2}x &= \frac{2}{3}y \\ \Leftrightarrow 3x &= 4y \\ \Leftrightarrow 3x - 4y &= 0 \quad (2) \end{align*}
		
		Từ (1) và (2), ta có hệ phương trình:
		$$ \begin{cases} x + y = 245 \\ 3x - 4y = 0 \end{cases} $$
		Giải hệ phương trình:
		\begin{align*} \begin{cases} x = 140 \\ y = 105 \end{cases} \end{align*}
		(Thỏa mãn điều kiện $x, y \in \mathbb{N}^*$).
		
		Kiểm tra lại:
		Số sách Toán đã phát: $\frac{1}{2} \cdot 140 = 70$ quyển.
		Số sách Ngữ văn đã phát: $\frac{2}{3} \cdot 105 = 2 \cdot 35 = 70$ quyển.
		(Số sách Toán đã phát bằng số sách Ngữ văn đã phát, hợp lý).
		Tổng số sách: $140 + 105 = 245$ quyển (hợp lý).
		
		Vậy, Hội khuyến học tỉnh đã tặng cho trường A 140 quyển sách Toán và 105 quyển sách Ngữ văn.
	}
\end{ex}
\begin{ex}
	Cho hình vuông ABCD có hai đường chéo cắt nhau tại O. Trên các cạnh AB, BC lần lượt lấy các điểm E, F sao cho $\widehat{EOF} = 90^\circ$.
	\begin{enumerate}
		\item[a)] Chứng minh tứ giác EBFO là tứ giác nội tiếp và tính $\widehat{EFO}$.
		\item[b)] Tia AF cắt tia DC tại N, tia OF cắt BN tại M. Chứng minh EF//BN và $FO \cdot FM = FB \cdot FC$.
	\end{enumerate}
\end{ex}

